\documentclass[a5paper,10pt]{article}

% https://github.com/InDevRus/filippov-solutions

% Нумерация страниц
\pagenumbering{gobble}

% Настройка полей
\usepackage{geometry}
\geometry{left=1cm}
\geometry{right=1cm}
\geometry{top=1cm}
\geometry{bottom=1.5cm}

% Вставка таблицы
\usepackage{array}
\usepackage{tabularx}

% Инструменты выравнивания формул
\usepackage{amsmath}
\usepackage{mathtools}

% Диакритика в математике
\usepackage{amsfonts}

% Вычисления размеров на ходу
\usepackage{calc}

% Удобные записи производных
\usepackage[italic=true]{derivative}

% Настройки сносок
\usepackage{enotez}

% Длинные знаки векторов
\usepackage{anyfontsize}
\usepackage[b]{esvect}

% Вставка картинок
\usepackage{graphicx}

% Вставка красной строки
\usepackage{indentfirst}
\setlength{\parindent}{1em}

% Условия в командах
\usepackage{xifthen}

% Луа-скрипты
\usepackage{luacode}

% Настройка команд отдельно в математическом и текстовом режимах
\usepackage{mathcommand}

% Для повторения LaTeX кода
\usepackage{multido}

% Конфигурация отступов формул
\delimitershortfall-1sp
\usepackage{mleftright}
\mleftright

% Растягивание объектов
\usepackage{relsize}

% Обтекание текстом
\usepackage{wrapfig}
\usepackage{subcaption}
\usepackage{floatflt}

% Поддержка ввода кириллицы
\usepackage[english, russian]{babel}

% Настройка корневой позиции для компилятора
\usepackage{import}
\providecommand*{\ProjectRootPath}{..}

\import{\ProjectRootPath/preambles/macros/}{theorems}

\import{\ProjectRootPath/preambles/fonts}{font_setup}

% Знак градуса
\usepackage{siunitx}

\import{\ProjectRootPath/preambles/macros/}{operators}
\import{\ProjectRootPath/preambles/macros/}{redefinitions}

\import{\ProjectRootPath/preambles/fonts}{letters_setup}
\import{\ProjectRootPath/preambles/fonts/adjustments}{custom_superscripts}

\import{\ProjectRootPath/preambles/custom}{formatting}
\import{\ProjectRootPath/preambles/custom}{frames}
\import{\ProjectRootPath/preambles/custom}{list_setup}

% TODO: перевести команды на современный стиль объявления

\begin{document}

Для касания окружности оси абсцисс необходимо и достаточно, чтобы её радиус равнялся модулю ординаты центра. Таким образом, вид уравнения такой окружности таков: $\br{x - a}^2 + \br{y - b}^2 = \pow{b}{2}$.
$\br{x - a}^2 + \pow{y}{2} - 2by = 0$.
$\br{x - a} + yy\` - by\` = 0$.
Снова дифференцируем: $1 + \br{\y\`}^2 + y\primesub[2]{y}{} - b\primesub[2]{y}{} = 0$.
Наконец, подставляем в исходное уравнение 
$\tsys{& b = \dfrac {1 + \br{\y\`}^2 + y\primesub[2]{y}{}} {\primesub[2]{y}{}} \\& x - a = by\` - yy\`}$,
откуда 
$x - a = \dfrac {1 + \br{\y\`}^2 + y\primesub[2]{y}{}} {\primesub[2]{y}{}} y\` - yy\` = \dfrac {y\` + \br{\y\`}^3} {\primesub[2]{y}{}}$. \\
$\br{\dfrac {y\` + \br{\y\`}^3} {\primesub[2]{y}{}}}^2 + \pow{y}{2} - 2y \dfrac {1 + \br{\y\`}^2 + y\primesub[2]{y}{}} {\primesub[2]{y}{}} = 0$. \\
$\br{y\` + \br{\y\`}^3}^2 - \pow{y}{2}\br{\primesub[2]{y}{}}^2 - 2y\primesub[2]{y}{} - 2y\br{\y\`}^2\primesub[2]{y}{} = 0$. \\
$\br{\y\`}^2 + 2\br{\y\`}^4 + \br{\y\`}^6 - \pow{y}{2}\br{\primesub[2]{y}{}}^2 - 2y\primesub[2]{y}{} - 2y\br{\y\`}^2\primesub[2]{y}{} = 0$. \\
$\br{\br{\y\`}^6 + 3\br{\y\`}^4 + 3\br{\y\`}^2 + 1} - \br{\br{\y\`}^4 + \pow{y}{2}\br{\primesub[2]{y}{}}^2 + 1 + 2\br{\y\`}^2 + 2y\primesub[2]{y}{} + 2y\br{\y\`}^2\primesub[2]{y}{}} = 0$. \\
$\br{\br{\y\`}^2 + 1}^3 - \br{\br{\y\`}^2 + y\primesub[2]{y}{} + 1}^2 = 0$.

\end{document}

